\documentclass[11pt]{article}
\usepackage[a4paper,margin=2.2cm]{geometry}
\usepackage[T1]{fontenc}
\usepackage{lmodern}
\usepackage{microtype}
\usepackage{amsmath}
\usepackage{graphicx}
\usepackage{booktabs}
\usepackage{adjustbox}
\usepackage{longtable}
\usepackage{float}
\usepackage{caption}
\usepackage{listings}
\usepackage[numbers,sort&compress]{natbib}
\usepackage[hidelinks]{hyperref}
\lstset{basicstyle=\scriptsize\ttfamily,breaklines=true,columns=fullflexible,keepspaces=true}
\captionsetup{font=small}
\input{numbers.tex}

\title{Adapting VERA to Firearm Categories and Actions by Verbalized Learning (Training 1)}
\author{VISILAB firearm-action dataset, frozen InternVL2-8B}
\date{\today}

\begin{document}
\maketitle

\begin{abstract}
VERA adapts a frozen vision-language model to a video task by letting the model rewrite the guiding questions it answers before it decides. We applied VERA's training procedure unchanged to the VISILAB firearm-action dataset with a frozen InternVL2-8B. The model explains the scene through its answers to the questions, then names the category (Handgun, Machine\_Gun or No\_Gun), the action group and the action code. Training used \val{split_train_videos} videos for 1,200 steps, with selection on \val{split_val_videos} validation videos and a single test run on \val{split_test_videos} videos. In the first run, the optimizer rewrote the questions in \val{opt_updated} of \val{opt_steps} steps. No learned checkpoint beat the neutral seed questions on validation: the best scored \val{val_selected_category_acc_pct}\,\% category accuracy against \val{val_seed_category_acc_pct}\,\%. On the test split the selected questions reached \val{test_learned_category_acc_pct}\,\% category accuracy against \val{test_seed_category_acc_pct}\,\% for the seed questions. The paired difference was \val{test_learnedminusseed_category_acc_pct} points, with a 95\,\% interval of \val{test_learnedminusseed_category_low_pct} to \val{test_learnedminusseed_category_high_pct}, so no improvement is shown. Action-group and exact-code accuracy did not improve detectably either. The selected questions ask whether a firearm is held and of what kind, and they recognise more handguns, but most handgun videos are still called No\_Gun. The run took about six GPU hours on one RTX~6000 Ada.\IfFileExists{./variants_section.tex}{ Two more training orders gave \val{cmp_vera_20261027_test_category_pct} and \val{cmp_vera_20261028_test_category_pct}\,\% test accuracy. Two declared variants were also run: a fixed loop with score memory and accept-if-not-worse, and a confusion-driven yes/no question bank. Neither reliably beat VERA or the seed questions; the fixed variant's three runs scored \val{cmp_vera_fixed_20261026_test_category_pct}, \val{cmp_vera_fixed_20261027_test_category_pct} and \val{cmp_vera_fixed_20261028_test_category_pct}\,\%. The spread across training orders was as large as any difference between methods.}{}

\end{abstract}

\section{Aim}
VERA~\cite{ye2025vera} adapts a frozen vision-language model to a new video task without changing its weights: the same model, acting as an optimizer, rewrites the guiding questions that the learner answers before it gives its decision, and the questions with the best validation score are kept. Training~1 applies this procedure to the VISILAB firearm-action dataset~\cite{ruizsantaquiteria2024firearm}. Each answer explains what the model sees (its answers to the guiding questions, as in VERA on UCF-Crime~\cite{sultani2018ucfcrime}) and then names the category (Handgun, Machine\_Gun, No\_Gun), the action group and the exact action code. Boxes are not used; they are Training~2.

\section{Data and split}
All 398 videos are used. The category is the dataset folder; the action code is the folder name's first field (41 codes; PAW7 of the paper has no video). Codes map to six action groups: carry and aim for firearm videos (reload variants have aim codes) and move, phone, bottle and heavy object for No\_Gun videos; descriptions are composed from the paper's Tables~1--3. Repeated takes of one scripted action by one subject at one place (a recording group) are near-duplicates and never span two splits. Per action code, recording groups were shuffled with a fixed seed and assigned so that every code appears in every split, with video targets of 60/15/25\,\%: \val{split_train_videos} training, \val{split_val_videos} validation and \val{split_test_videos} test videos (Table~\ref{tab:S1}). A subject-held-out split is impossible because the handgun videos were filmed with two of the four subjects only.

\input{tables/S1.tex}

\section{Method}
\paragraph{Faithful to VERA's \texttt{training.py}.} Learner and optimizer are the same frozen InternVL2-8B~\cite{chen2024internvl15} (revision 6fb9ad69, CUDA BF16, FlashAttention, greedy decoding, one beam, 1,024 new tokens). A batch holds two training videos with eight frames each, sampled evenly over the whole video as VERA's training data loader does. The learner answers both videos in one \texttt{batch\_chat} call; the optimizer sees both videos' sixteen frames, the current questions, the batch's predictions and targets, and answers with reasoning followed by ``New Prompt Questions'', from which VERA's rule takes at most five lines. The prompt assembly and the question update were checked against VERA's own code by executing its \texttt{forward} and \texttt{training\_step} on a fake model. Validation runs every 100 training batches within an epoch over ten epochs; the validated questions with the highest validation category accuracy (then action-group accuracy, then the earliest step) are selected.

\paragraph{Adapted to this task.} VERA's binary model description and its optimizer instruction were rewritten for three categories, six action groups and 41 codes, with the same structure and wording elsewhere (Appendix~\ref{app:prompts}). Predictions and targets given to the optimizer are [category, action group, code] triples. Answers are parsed strictly: an unparseable field is invalid and counts as wrong, never as a default. The seed questions are neutral for the task. The training order is shuffled with a fixed seed. The seed questions are also validated before training, for reference only. VERA's temporal score refinement is not used: it is defined for one binary anomaly score.

\paragraph{Evaluation.} The seed questions and the selected questions are each run once on the test split. Accuracy is video-level; intervals are 95\,\% percentile intervals from 2,000 bootstrap draws that resample recording groups within category (seed 20261026); the learned-minus-seed difference is paired on the same videos.

\section{Results}
% PROSE PENDING: written from the saved analysis after the run. Floats below are final.
\subsection{Training}
\input{tables/S2.tex}
\begin{figure}[htbp]\centering\includegraphics[width=\textwidth]{figures/F1_validation_trajectory.pdf}
\caption{Validation accuracy of the learned questions at each validation step; dashed lines are the seed questions.}\label{fig:F1}\end{figure}
\input{tables/S9.tex}

\subsection{Test}
\input{tables/S4.tex}
\input{tables/S5.tex}
\begin{figure}[htbp]\centering\includegraphics[width=\textwidth]{figures/F2_category_confusion.pdf}
\caption{Category confusion on the test split for the seed and the selected questions.}\label{fig:F2}\end{figure}
\input{tables/S6.tex}
\begin{figure}[htbp]\centering\includegraphics[width=\textwidth]{figures/F3_accuracy_by_category.pdf}
\caption{Test accuracy by true category, seed versus selected questions.}\label{fig:F3}\end{figure}
\input{tables/S7.tex}
\input{tables/S8.tex}

\subsection{Computation}
\input{tables/S10.tex}
\begin{figure}[htbp]\centering\includegraphics[width=\textwidth]{figures/F4_runtime.pdf}
\caption{Learner and optimizer time per training step.}\label{fig:F4}\end{figure}

Appendix~\ref{app:prompts} lists the prompts and the questions; Appendix~\ref{app:cases} shows answers chosen by a fixed rule.


\section{Limitations}
\begin{itemize}
\item \textbf{Small validation and test sets.} Selection uses \val{split_val_videos} validation videos, so one video moves validation accuracy by 1.6 points, and the selected checkpoint can be a lucky one. The test split has \val{split_test_videos} videos in few recording groups per action code, so test intervals are wide and exact-code accuracy is exploratory.
\item \textbf{One split, one model, one decoding.} The results hold for this split, the frozen InternVL2-8B at one revision, eight frames per video and greedy decoding. A subject-held-out split is impossible: the handgun videos show two of the four subjects.
\item \textbf{The seed questions and the prompts are ours.} VERA's prompts were rewritten for three categories and six action groups, and the neutral seed questions are our choice. A different start or wording may train differently.
\item \textbf{Strict parsing.} An unparseable field counts as wrong, so accuracy measures answering in the required format as well as recognising the scene. Greedy repetition loops that reach the 1,024-token cap are counted, not suppressed.
\item \textbf{Faithful means faithful to VERA's training loop.} The optimizer keeps no record of earlier scores and rewrites the questions after every batch, as in VERA.\IfFileExists{./variants_section.tex}{ The variants in Section~\ref{sec:variants} change this; the faithful run does not.}{}
\IfFileExists{./variants_section.tex}{\item \textbf{Three runs per method, one for each bank variant.} The spread across run seeds is estimated from three runs of VERA and three of the fixed variant; the bank and its ablations ran once. The variants were designed after seeing the first faithful run's training logs, though not its test results.}{}
\item \textbf{No boxes, no temporal refinement.} Box supervision is Training~2. VERA's temporal score refinement is defined for one binary anomaly score and is not used.
\end{itemize}


\bibliographystyle{unsrtnat}
\bibliography{references}

\appendix
\section{Prompts and questions}\label{app:prompts}
\subsection*{Learner instruction}
\lstinputlisting{learner_instruct.txt}
\subsection*{Optimizer instruction}
\lstinputlisting{optimizer_instruct.txt}
\subsection*{Seed questions}
\lstinputlisting{questions_seed.txt}
\subsection*{Selected questions}
\lstinputlisting{questions_selected.txt}

\section{Qualitative answers}\label{app:cases}
For each true category, the first test video (by ID) whose category the selected questions got right and the first they got wrong.
\input{qualitative.tex}
\end{document}
